\section{Introduction}
In a conventional battery-electric drive, the battery is a single high-voltage string that feeds a two-level (2L) voltage-source inverter. With silicon carbide (SiC) devices, such inverters reach peak efficiencies above \SI{99}{\percent}, but they switch the full DC-link voltage at a fixed carrier frequency of typically \SI{10}{\kilo\hertz} or more, and they need a bulky DC-link capacitor.

The battery-integrated modular multilevel converter (BI-MMC), also known as an AC battery or modular multilevel battery \cite{lesnicar2003,helling2016,quraan2016}, takes a different route. It distributes the cells over $\Nsm$ series-connected full-bridge submodules (SMs) per phase. Each SM inserts, bypasses or reverses its local cell string, so the phase voltage is synthesised directly from the batteries. This permits low-voltage Si MOSFETs with very low on-resistance, SM switching at only a few times the fundamental frequency, and inherent cell balancing. The price is that the cells carry the switched phase current, including its reactive and field-weakening components, which raises the battery Joule losses \cite{kersten2019}.

Whether the BI-MMC is the more efficient drive therefore depends on where in the torque--speed plane a vehicle spends its energy. Published comparisons usually consider a few operating points or a single vehicle, or they neglect the battery. A fair comparison needs (a)~the same machine, cells and installed energy in both drives, (b)~the battery losses of both drives, and (c)~a weighting of the operating points by a realistic duty cycle.

\paragraph{Contributions.} This paper documents and applies an open-source evaluation platform that provides exactly this. It makes three contributions:
\begin{enumerate}
\item \emph{Efficiency maps.} The platform computes the efficiency map of the 2L inverter and of the BI-MMC over the full torque--speed plane for motoring and regeneration, with identical machine, cells and energy, and with MTPA currents solved separately for each drive's voltage limit (Analysis~I, \cref{sec:maps}).
\item \emph{Drive-cycle efficiency.} The same loss equations are evaluated along a duty cycle, sample by sample, with the state of charge (SOC) and the device, case and cell temperatures tracked, and the cycle efficiency of both drives is returned (Analysis~II, \cref{sec:cycles}).
\item \emph{Complete documentation and a study.} All equations, algorithms, data and assumptions of the Python implementation are written down (\cref{sec:software,sec:models} and the appendices), and the platform is applied to six applications from passenger cars to trains (\cref{sec:results}). The study yields a simple rule for when an AC battery pays off.
\end{enumerate}

\paragraph{Organisation.} \Cref{sec:system} defines the compared systems and the notation. \Cref{sec:software} describes the software architecture. \Cref{sec:models} gives the models shared by both analyses: machine design, battery sizing, MTPA control, and the loss models of the two drives. \Cref{sec:maps,sec:cycles} describe the two analyses. \Cref{sec:results} presents the results for six applications, and \cref{sec:discussion} discusses assumptions and limitations. The appendices list all parameters, a function reference, the output formats and per-application traces.
